% ============================================================
% §6 模型评价与推广 (Stage 07 产物) — 精简版
% 竞赛: 电工杯 B题 (2026)
% 详细版见附录E
% ============================================================

\section{模型评价与推广}

\subsection{模型优点}

\textbf{优点1：随机种子揭示24\%目标值提升。} 通过设定randomSeed=42的确定性求解，CBC发现了目标值54.225的全覆盖最优解（vs非确定解43.644），揭示CFLP类MILP中多起点随机重启的方法论价值。

\textbf{优点2：MILP-MCLP-ES精确线性化框架。} 针对$S_2(u_i)\cdot x_{ij}$非线性乘积，采用McCormick包络以4条线性约束等价替换；$S_1$（4段）与$S_2$（5段）各以Big-M方法分段线性化，将原始NLP等价转化为MILP，保留全局最优性保证。

\textbf{优点3：多维鲁棒性超越强稳健性阈值。} Q4三场景7组参数配置下，覆盖率、满意度、利润率韧性系数均大于0.80，远超0.75的强鲁棒性阈值。预算从120万增至130万（+8.3\%）触发覆盖相变，为政策制定提供了精确的边际投资锚点。

\textbf{优点4：模型输出具备直接工程可操作性。} Q2选址方案可直接转化为招标标段划分，Q3定价已核算至0.01元精度，Q4预算相变分析为"追加多少投资可实现全覆盖"提供了定量答案。

\subsection{模型局限与改进方向}

\textbf{局限1：Q2 MILP对随机种子的敏感性}——多起点随机重启（10次独立求解取最优）可将期望目标值提升3--8\%，代价为计算时间$\times$10。

\textbf{局限2：S2全员触底的满意度悬崖}——引入软容量约束或85\%缓冲系数可释放0.600$\to$0.720的S2跃迁，但需追加13万元容量补偿。

\textbf{局限3：Markov转移概率的平稳性假设}——采用时变转移矩阵$\mathbf{A}(t)$或贝叶斯层次模型（PyMC/MCMC后验）可捕捉队列效应的时变性，代价为需3--5年纵向追踪数据。

\textbf{局限4：序贯求解（Q1$\to$Q2$\to$Q3）非联合最优}——Benders分解可实现联合优化（预期提升3--5\%），但开发时间+15--20h对竞赛窗口不可行，更适合赛后学术扩展。

\textbf{改进方向}：(1) 分布鲁棒优化（DRO）以Wasserstein模糊集替代确定性需求，CVXPY 1.2+ SOCP求解；(2) 两阶段随机规划引入"观望-决策"灵活性，SAA方法求解，期望后悔值降低15--20\%。

\subsection{模型推广}

本框架的MCLP+满意度+三层交叉补贴核心方法论可迁移至三类公共设施优化场景：
\textbf{(1) 分级诊疗站点选址}——"15分钟医疗服务圈"与MCLP覆盖目标同构，直接对接国家卫健委2027年75\%签约覆盖率政策目标；
\textbf{(2) 冷链前置仓网络规划}——三类老人$\to$三温区商品品类，配送时效满意度替代距离满意度，前置仓亏损率30--40\%的行业痛点可通过选址-容量-定价联合优化缓解；
\textbf{(3) 新能源充电桩布局}——消除"充电荒漠"与MCLP全覆盖目标精确对齐，补贴帽机制可防止财政资源错配。
各场景的映射关系与调整假设已在上文详述。
